\item El ingeniero de planeación de una nueva planta de alumbre debe presentar a su compañía un cálculo aproximado de la capacidad de un silo diseñado para contener mineral de bauxita hasta que sea procesado en alumbre. El mineral tiene la apariencia de talco rosado y se vierte desde una banda transportadora situada en la parte superior del silo. El silo es un cilindro de 100 pies de altura con un radio de 200 pies. La banda transporta $60\ 000 \, \pi \text{ pies}^3/\text{h}$ y el mineral se apila manteniendo una forma cónica cuyo radio es 1.5 veces su altura.
\begin{enumerate}
    \item[(a)] Si en cierto tiempo $t$ la pila tiene 60 pies de altura, ¿cuánto tiempo transcurrirá para que la pila alcance la parte superior del silo?
    \item[(b)] La administración desea saber cuánto espacio quedará libre en el área del piso del silo cuando la pila tenga 60 pies de altura. ¿Con qué rapidez crece el área del piso de la pila cuando alcanza esa altura?
    \item[(c)] Supóngase que una máquina cargadora empieza a sacar el mineral a razón de $20\ 000 \, \pi \text{ pies}^3/\text{h}$, cuando la altura de la pila alcanza los 90 pies. Supóngase también que la pila continúa manteniendo su forma cónica. ¿Cuánto tiempo transcurrirá para que bajo estas condiciones la pila alcance la parte superior del silo?
\end{enumerate}
